\chapter{Cha Mẹ Đi Qua Cửa Lớp}
\markboth{Cha Mẹ Đi Qua Cửa Lớp}{Parents Walk Through the Classroom Door Too}

\section{Kỳ Vọng Đi Trước Đứa Trẻ}
\begin{truyen}{Kỳ Vọng Đi Trước Đứa Trẻ}{Expectations Arrive Before the Child}
\chuhoa{C}{ó phụ huynh chưa cần nghe con nói đã biết sẵn con phải đứng ở đâu.}
\textit{(Some parents already know where the child must stand before hearing the child speak.)}

Kỳ vọng mạnh có thể giúp một đứa trẻ cố gắng.
\textit{(Strong expectation can push a child to try harder.)}

Nhưng nếu đi quá xa trước sức của con, nó biến thành một bóng đè rất đẹp bề ngoài.
\textit{(But if it runs too far ahead of the child’s capacity, it becomes a beautifully dressed pressure.)}

Không phải đứa trẻ nào cũng ngã vì lười. Có em ngã vì phải chạy bằng chân của người khác.
\textit{(Not every child falls because of laziness. Some fall because they are running on someone else’s legs.)}

\end{truyen}

\begin{baihoc}
Kỳ vọng tốt nhất là kỳ vọng biết đi cùng nhịp lớn lên của con.
\textit{(The best expectation knows how to walk with the child’s growth pace.)}
\end{baihoc}

\begin{ghinhoanh}
\textbf{Short English to remember:}

The best expectation walks with the child’s pace of growth.

\end{ghinhoanh}

\begin{tuhocnhanh}
\textbf{Gợi ý để dạy và tự nghĩ / Teaching and reflection prompts}

1. Kỳ vọng này đang đỡ hay đè? \textit{Is this expectation supporting or crushing?}

2. Vì sao phải đi cùng nhịp của con? \textit{Why must expectation walk at the child’s pace?}

3. Người dạy nói với phụ huynh điều này thế nào? \textit{How can a teacher say this to parents?}
\end{tuhocnhanh}
\ngancach

\section{Bênh Con Quá Mức Cũng Làm Hại Con}
\begin{truyen}{Bênh Con Quá Mức Cũng Làm Hại Con}{Overprotecting a Child Can Also Harm the Child}
\chuhoa{C}{ó phụ huynh thương con thật, nên hễ con buồn là muốn tìm ngay người khác sai.}
\textit{(Some parents truly love their child, so whenever the child is upset they immediately search for someone else to blame.)}

Sự bênh vực đó giúp con đỡ đau nhất thời.
\textit{(That protection eases pain in the moment.)}

Nhưng nếu kéo dài, nó tước khỏi con cơ hội học trách nhiệm và đối diện phần sai của mình.
\textit{(But over time, it steals the child’s chance to learn responsibility and face personal mistakes.)}

Yêu con không phải lúc nào cũng là che cho con khỏi mọi va chạm.
\textit{(Loving a child does not always mean shielding the child from every collision.)}

\end{truyen}

\begin{baihoc}
Có những lần để con tự nhìn phần sai của mình mới là thương đúng cách.
\textit{(Sometimes the right kind of love lets a child see their own mistake clearly.)}
\end{baihoc}

\begin{ghinhoanh}
\textbf{Short English to remember:}

Sometimes love means letting a child face their own mistake.

\end{ghinhoanh}

\begin{tuhocnhanh}
\textbf{Gợi ý để dạy và tự nghĩ / Teaching and reflection prompts}

1. Sự bênh này đang cứu hay làm yếu con? \textit{Is this protection helping or weakening the child?}

2. Vì sao đối diện sai là cần thiết? \textit{Why is facing mistakes necessary?}

3. Người dạy có thể trao đổi ra sao để không gây đối đầu? \textit{How can a teacher discuss this without creating conflict?}
\end{tuhocnhanh}
\ngancach

\section{Cha Mẹ Cũng Đang Quá Mệt}
\begin{truyen}{Cha Mẹ Cũng Đang Quá Mệt}{Parents May Also Be Carrying Too Much}
\chuhoa{N}{hiều phụ huynh đến trường với vẻ khó chịu, gấp gáp, hoặc lạnh. Không phải vì họ không thương con. Có khi vì họ đã mệt quá lâu.}
\textit{(Many parents arrive at school tense, rushed, or cold. Not because they do not love their child, but because they have been tired for too long.)}

Tiền bạc, công việc, bệnh tật, và áp lực lớn làm cách họ nói chuyện trở nên cứng hơn.
\textit{(Money, work, illness, and heavy pressure can harden the way they speak.)}

Nhìn đúng, người dạy sẽ bớt giận mà hiểu hơn.
\textit{(Seen rightly, teachers may feel less angry and more understanding.)}

Lớp học không chỉ có học sinh mang gánh nặng. Người lớn đi cùng các em cũng có phần nặng của riêng họ.
\textit{(Students are not the only ones carrying burdens. The adults around them do too.)}

\end{truyen}

\begin{baihoc}
Muốn làm việc với gia đình tốt hơn, đôi khi phải nhìn phụ huynh như những người cũng đang cố sống qua ngày.
\textit{(To work better with families, we sometimes must see parents as people also trying to survive their days.)}
\end{baihoc}

\begin{ghinhoanh}
\textbf{Short English to remember:}

Parents are also people trying to survive their days.

\end{ghinhoanh}

\begin{tuhocnhanh}
\textbf{Gợi ý để dạy và tự nghĩ / Teaching and reflection prompts}

1. Điều gì có thể đang đè lên phụ huynh? \textit{What may be weighing on the parent?}

2. Vì sao góc nhìn này giúp người dạy mềm hơn? \textit{Why does this perspective help teachers soften?}

3. Sự hợp tác với gia đình cần bắt đầu bằng điều gì? \textit{What should collaboration with families begin with?}
\end{tuhocnhanh}
\ngancach

\section{Một Cuộc Nói Chuyện Có Thể Đổi Hướng Một Đứa Trẻ}
\begin{truyen}{Một Cuộc Nói Chuyện Có Thể Đổi Hướng Một Đứa Trẻ}{One Conversation Can Shift a Child’s Direction}
\chuhoa{C}{ó khi cả năm học căng thẳng, nhưng chỉ một buổi ngồi lại đúng cách giữa thầy cô và cha mẹ là mọi thứ dịu hẳn.}
\textit{(Sometimes an entire tense school year softens after one proper conversation between teacher and parent.)}

Khi người lớn thôi hơn thua và cùng nhìn về phía đứa trẻ, lớp học bớt chiến trường đi rất nhiều.
\textit{(When adults stop competing and turn together toward the child, the classroom stops feeling like a battlefield.)}

Điều cứu một học sinh nhiều khi không phải là một biện pháp mới, mà là người lớn chịu nghe nhau thật.
\textit{(What saves a student is often not a new strategy, but adults truly listening to one another.)}

Có những đứa trẻ đổi hướng từ việc người lớn cuối cùng cũng nói cùng một tiếng nói.
\textit{(Some children change direction when the adults around them finally speak with one voice.)}

\end{truyen}

\begin{baihoc}
Khi người lớn hợp lực đúng cách, học sinh bớt phải chịu những vết kéo từ nhiều phía.
\textit{(When adults work together well, students suffer less from being pulled in many directions.)}
\end{baihoc}

\begin{ghinhoanh}
\textbf{Short English to remember:}

When adults work together well, students suffer less from conflicting pulls.

\end{ghinhoanh}

\begin{tuhocnhanh}
\textbf{Gợi ý để dạy và tự nghĩ / Teaching and reflection prompts}

1. Điều gì làm một cuộc nói chuyện trở nên có ích? \textit{What makes a conversation truly useful?}

2. Vì sao học sinh khổ khi người lớn kéo nhiều phía? \textit{Why do students suffer when adults pull in different directions?}

3. Người dạy có thể chủ động phần nào? \textit{What part can the teacher take initiative in?}
\end{tuhocnhanh}
\ngancach

\section{Có Phụ Huynh Nói Rất To Vì Đang Quá Lo}
\begin{truyen}{Có Phụ Huynh Nói Rất To Vì Đang Quá Lo}{Some Parents Speak Loudly Because They Are Too Afraid}
\chuhoa{K}{hông phải phụ huynh nào nói gắt cũng chỉ vì muốn hơn thua. Có người đang sợ con mình tụt lại, sợ mình bất lực, sợ không biết phải làm gì.}
\textit{(Not every sharp parent speaks from ego alone. Some are afraid their child is falling behind, afraid of being helpless, afraid because they do not know what to do.)}

Lo lắng nếu đi qua thiếu bình tĩnh sẽ thành áp lực.
\textit{(Anxiety, when passing through a lack of calm, turns into pressure.)}

Người dạy cần nghe được nỗi sợ đằng sau giọng nói.
\textit{(Teachers need to hear the fear behind the tone.)}

Hiểu điều đó không phải để nhường sai. Mà để nói đúng hơn.
\textit{(Understanding this does not mean yielding to what is wrong. It means responding more accurately.)}

\end{truyen}

\begin{baihoc}
Đằng sau một giọng nói căng có khi là một nỗi lo chưa biết cách nói cho tử tế.
\textit{(Behind a tense voice there may be a fear that has not learned how to speak kindly.)}
\end{baihoc}

\begin{ghinhoanh}
\textbf{Short English to remember:}

Behind a tense voice there may be a fear that has not learned how to speak kindly.

\end{ghinhoanh}

\begin{tuhocnhanh}
\textbf{Gợi ý để dạy và tự nghĩ / Teaching and reflection prompts}

1. Nỗi lo nào có thể đứng sau giọng nói này? \textit{What fear may stand behind this tone?}

2. Hiểu nỗi sợ khác với nhượng bộ ở đâu? \textit{How is understanding different from giving in?}

3. Người dạy có thể đáp lại thế nào cho đúng hơn? \textit{How can a teacher respond more accurately?}
\end{tuhocnhanh}
\ngancach

\section{Có Phụ Huynh Gặp Mình Chỉ Để Xin Đừng Bỏ Con Họ}
\begin{truyen}{Có Phụ Huynh Gặp Mình Chỉ Để Xin Đừng Bỏ Con Họ}{Some Parents Only Ask One Thing: Please Do Not Give Up on My Child}
\chuhoa{C}{ó những cuộc gặp rất ngắn. Phụ huynh không biện minh nhiều. Chỉ nói: “Mong thầy cô để ý giùm cháu. Ở nhà chúng tôi nói không nổi nữa.”}
\textit{(Some meetings are very short. Parents do not defend much. They only say, “Please keep an eye on my child. We cannot reach them at home anymore.”)}

Câu đó nghe mệt, nhưng cũng nghe rất thật.
\textit{(That sentence sounds heavy, but very real.)}

Đó là lúc người lớn bên này nhìn thấy người lớn bên kia cũng đang đuối.
\textit{(It is the moment one adult sees that the other adult is exhausted too.)}

Giáo dục nhiều khi không bắt đầu bằng kế hoạch lớn. Nó bắt đầu bằng việc chưa ai bỏ em đó trước.
\textit{(Education often does not begin with a big plan. It begins with no adult giving up on that child first.)}

\end{truyen}

\begin{baihoc}
Giữ một học sinh lại nhiều khi là việc người lớn giữ nhau cùng đứng lại trước đã.
\textit{(Holding onto a student often begins with adults holding their place together first.)}
\end{baihoc}

\begin{ghinhoanh}
\textbf{Short English to remember:}

Sometimes helping a student begins with adults refusing to give up first.

\end{ghinhoanh}

\begin{tuhocnhanh}
\textbf{Gợi ý để dạy và tự nghĩ / Teaching and reflection prompts}

1. Câu xin giúp này cho thấy điều gì? \textit{What does this plea reveal?}

2. Vì sao người lớn cần đứng cùng nhau? \textit{Why must adults stand together?}

3. Mình có từng vô tình bỏ một em quá sớm không? \textit{Have we ever given up on a student too early?}
\end{tuhocnhanh}
\ngancach

\section{Đừng Chỉ Gọi Phụ Huynh Khi Có Chuyện}
\begin{truyen}{Đừng Chỉ Gọi Phụ Huynh Khi Có Chuyện}{Do Not Contact Parents Only When Something Is Wrong}
\chuhoa{N}{ếu số điện thoại của giáo viên chỉ hiện lên lúc học sinh có vấn đề, thì phụ huynh sẽ quen với tâm thế sợ.}
\textit{(If a teacher’s number appears only when there is trouble, parents begin to live in a state of fear.)}

Một lời báo tốt, một lời khen đúng lúc, một cuộc gọi ngắn để nói điều tích cực cũng rất quan trọng.
\textit{(A positive note, a timely compliment, a short call carrying good news also matter greatly.)}

Khi mối liên hệ không chỉ toàn cảnh báo, sự hợp tác sẽ dễ thở hơn.
\textit{(When communication is not made only of warnings, cooperation becomes easier to breathe in.)}

Người lớn cũng cần những cây cầu tốt đẹp, không chỉ cần những chuông báo động.
\textit{(Adults also need bridges of goodwill, not only alarm bells.)}

\end{truyen}

\begin{baihoc}
Quan hệ giữa gia đình và nhà trường bền hơn khi nó có cả tin vui, không chỉ tin xấu.
\textit{(The relationship between home and school lasts better when it carries good news as well as bad.)}
\end{baihoc}

\begin{ghinhoanh}
\textbf{Short English to remember:}

Partnership grows better when it carries good news too.

\end{ghinhoanh}

\begin{tuhocnhanh}
\textbf{Gợi ý để dạy và tự nghĩ / Teaching and reflection prompts}

1. Vì sao chỉ gọi khi có chuyện là chưa đủ? \textit{Why is contacting parents only in trouble not enough?}

2. Một tin tốt ngắn có tác dụng gì? \textit{What can a short positive message do?}

3. Mình đã tạo được cầu nối tích cực nào chưa? \textit{Have we created any positive bridge yet?}
\end{tuhocnhanh}
\ngancach

\section{Mỗi Lời Nói Về Con Trước Mặt Phụ Huynh Đều Có Trọng Lượng}
\begin{truyen}{Mỗi Lời Nói Về Con Trước Mặt Phụ Huynh Đều Có Trọng Lượng}{Every Word About a Child Before the Parent Has Weight}
\chuhoa{K}{hi thầy cô nói về một học sinh trước mặt cha mẹ, câu chữ không chỉ đi vào tai người lớn. Nó còn có thể đi thẳng vào cách người đó nhìn con mình.}
\textit{(When teachers speak about a student before the parents, the words do not only enter adult ears. They can go straight into the way that parent sees the child.)}

Một câu thiếu công bằng có thể bị mang về nhà và lặp lại rất lâu.
\textit{(An unfair sentence can be carried home and repeated for a very long time.)}

Ngược lại, một lời nhìn đúng cũng có thể mở ra chỗ tin cho cả gia đình.
\textit{(On the other hand, a fair insight can open a place of trust for the whole family.)}

Vì vậy lời nói của người dạy luôn nên có trách nhiệm hơn một chút.
\textit{(That is why a teacher’s words should always carry extra responsibility.)}

\end{truyen}

\begin{baihoc}
Nói về một đứa trẻ là đang góp phần định hình cách người lớn khác đối xử với đứa trẻ đó.
\textit{(To speak about a child is to help shape how other adults will treat that child.)}
\end{baihoc}

\begin{ghinhoanh}
\textbf{Short English to remember:}

To speak about a child is to shape how others may treat that child.

\end{ghinhoanh}

\begin{tuhocnhanh}
\textbf{Gợi ý để dạy và tự nghĩ / Teaching and reflection prompts}

1. Vì sao lời nói trước phụ huynh có sức nặng? \textit{Why do words before parents carry such weight?}

2. Một câu thiếu công bằng có thể gây gì? \textit{What can an unfair sentence cause?}

3. Mình có đang nói về học sinh đủ công bằng chưa? \textit{Are we speaking about students fairly enough?}
\end{tuhocnhanh}
